\documentclass[10pt,a4paper]{amsart}
\usepackage[margin=19mm]{geometry}
\usepackage{amsmath,amssymb,mathtools}
\usepackage[scr=rsfs,cal=euler]{mathalfa}
\usepackage{xfrac}
\usepackage[unicode,hidelinks,bookmarks=false]{hyperref}
\newcommand{\ie}{i.e.\ }
\newcommand{\cf}{cf.\ }
\newcommand{\resp}{respectively\ }
\newcommand{\Subsection}{\subsection}
\providecommand{\nobreakdash}{\nobreak\hbox{-}}
\setlength{\emergencystretch}{2em}
\multlinegap0pt
\usepackage{amssymb}
\usepackage{bm}
\newtheorem{lemma}{Lemma}
\newtheorem{proposition}{Proposition}
\title{Global branches of Stokes waves of variable period on stratified fluids}
\author{Vladimir Kozlov$^a$$^b$}
\date{Source: arXiv:2606.19018v1 (CC BY 4.0)}
\begin{document}
\maketitle

\noindent\emph{Source and licence.} Global branches of Stokes waves of variable period on stratified fluids. Version arXiv:2606.19018v1 (CC BY 4.0), distributed by arXiv under the Creative Commons Attribution 4.0 International licence, http://creativecommons.org/licenses/by/4.0/, as recorded in the arXiv licence metadata for this exact version and shown on the abstract page https://arxiv.org/abs/2606.19018v1. Full licence notice: \texttt{licenses/arxiv-2606.19018-CC-BY-4.0.txt}.\par\smallskip

\noindent\textit{Editorial scope.} Counted unit: the article's proposition PM1a (source0.tex lines 541--545). The article's complete original proof of it is delivered with it (source0.tex lines 546--584, one \texttt{proof} environment of 39 lines). No mathematics is written, repaired or re-derived: the statement and the proof are the article's own lines, in the article's own notation, and no retained line is substituted or re-flowed.  Retained uncounted as the source-local context the proof needs: lines 176--178; lines 180--182; lines 183--185; lines 188--190; lines 206--218; lines 456--463.  Nothing was omitted from any retained span, so the package carries no omission record.  No retained line contains a citation, so the wrapper carries no bibliography and no bibliography entry.  No retained line contains a figure environment, an image-inclusion command or drawing code, so no image is delivered and the package declares no asset dependency.  Editorial formatting only, in addition: an AMS \texttt{amsart} preamble that restates the article's own declarations and macros used by the retained lines, namely the environments \texttt{lemma}, \texttt{proposition} on separate counters, together with the standard packages those lines need.  The retained environments print in this document as Lemma 1, Proposition 1 in order of appearance, whereas the article prints the same items with the numbering of its own full document.  The complete native source of the article as distributed in the arXiv e-print for this exact version is bundled unchanged as \texttt{sources/203110/source0.tex}.
\begin{equation}\label{J18aa}
\Psi^{''}(y)+\omega(y,\Psi(y))=0\;\;\mbox{on $(-d,0)$}
\end{equation}

\begin{equation}\label{J18ab}
\Psi(0)=0\;\;\mbox{and}\;\;\Psi(-d)=-p_0,
\end{equation}

\begin{equation}\label{J18ac}
\frac{1}{2}|\Psi'(0)|^2+g\rho(0)d=R.
\end{equation}

\begin{equation}\label{F11a}
|\rho(p)|+|\beta(p)|\leq C_1(1+|p|)\;\;\mbox{and}\;\;|\rho(p_2)-\rho(p_1)|+|\beta(p_2)-\beta(p_1)|\leq C_1|p_2-p_1|.
\end{equation}

We consider the following branch of the laminar solutions $(\Psi(y;s),d(s))$, which solve the problem
\begin{eqnarray}\label{F19a}
&&\Psi^{''}(y)+\omega(y,\Psi(y))=0\;\;\mbox{on $(-\infty,0]$},\nonumber\\
&&\Psi(0)=0,\;\;\Psi'(0)=s,
\end{eqnarray}
together with
\begin{equation}\label{A10a}
\Psi(d(s))=-p_0.
\end{equation}
Due to (\ref{F11a}) the problem (\ref{F19a}) is uniquely solvable and we denote its solution by $\Psi(y;s)$. We take the maximal interval $(-\infty,s_*)$ such that equation (\ref{A10a}) holds and $\Psi_y(d(s);s)<0$ for $s<s_*$. Clearly the pair $(\Psi(y;s),d(s))$ solves the problem (\ref{J18aa})-(\ref{J18ac}) with
\begin{equation}\label{A17a}
R={\mathcal R}(s):=\frac{1}{2}\Psi_y(0;s)^2+g\rho(0)d(s).
\end{equation}

\begin{lemma}\label{LemF11a} There exists $d_1>0$ such that the problem (\ref{J18aa}), (\ref{J18ab})
has a unique solution for each $d\leq d_1$. This solution satisfies
\begin{equation}\label{F11aa}
\Psi(y)=\frac{p_0y}{d}+w,\;\;|w|d^{-1}+|w'|\leq Cd\;\;\mbox{for $y\in [0,-d]$},
\end{equation}
where $C$ is independent of $y$.

\end{lemma}

\begin{proposition}\label{PM1a} Solutions to the problem (\ref{F19a}) satisfy the following monotonicity property
\begin{equation}\label{F7b}
\Psi(y,s_2)>\Psi(y,s_1)\;\;\mbox{for $y<0$ if $s_2<s_1$}.
\end{equation}
\end{proposition}

\begin{proof} Since $s_2<s_1$ the inequality (\ref{F7b}) is true for small $y$,
Assume that $\Psi(y,s_2)=\Psi(y,s_1)$ for a certain $y_0$, which is the largest negative $y$ with this property. Let $s_3=(s_2+s_1)/2$. Then there exists $y_1>y_0$ such that $\Psi(y_1,s_3)=\Psi(y_1,s_1)$ or $\Psi(y_1,s_3)=\Psi(y_1,s_2)$ and this is the largest such $y$. We put $\tilde{s}_1=s_2$, $\hat{s}_1=s_3$ or $\tilde{s}_1=s_3$, $\hat{s}_1=s_1$ depending on the above choice of $s_3$. Then we have $\Psi(y_1,\tilde{s}_1)=\Psi(y_1,\hat{s}_1)$ and  $\tilde{s}_1<\hat{s}_1$.

Continuing this procedure we can construct a sequences $\{\hat{s}_j\}$, $\{\tilde{s}_j\}$ and $\{y_j\}$ with the following properties
\begin{equation}\label{M7a}
s_2\leq\hat{s}_1\leq\hat{s}_2\cdots <\cdots \tilde{s}_2\leq\tilde{s}_1\leq s_1,\;\; y_1<y_2<\cdots <0
\end{equation}
and
 $$
 \Psi(y_j,\hat{s}_j)=\Psi(y_j,\tilde{s}_j)\;\;
 \mbox{and}\:\;\Psi(y,\hat{s}_j)<\Psi(y,\tilde{s}_j)\;\;\mbox{ for $y\in (y_j,0)$}.
 $$
Moreover
\begin{equation}\label{M7aa}
\tilde{s}_j-\hat{s}_j\to 0\;\;\mbox{ as $j\to\infty$}.
\end{equation}
We denote by $s_\dag$ the limit point of the sequences $\{\hat{s}_j\}$ and $\{\tilde{s}_1\}$, which is the same due to (\ref{M7a}) and (\ref{M7aa}). Let also $y_\dag$ the limit of the sequence $\{y_j\}$.

If $y_\dag=0$ then due to uniqueness in Lemma \ref{LemF11a} there are no two different solution coinciding for a small $|y|$. If $y_\dag>0$ then application of the assumption
(A) to the function $\Psi(\cdot,s_\dag)$ on the interval $(y_\dag,0)$ implies that the Dirichlet problem  for the Fr\'{e}chet derivative at the function $\Psi(\cdot,s_\dag)$ has positive first eigenvalue.

Furthermore, we have
\begin{eqnarray*}
&&\omega(y,\Psi(y;\hat{s}_j))-\omega(y,\Psi(y;\tilde{s}_j))=
\int_0^1\frac{d}{d\,t}\omega(y,t\Psi(y;\hat{s}_j)+(1-t)\Psi(y;\tilde{s}_j))dt\\
&&=\int_0^1\omega_p(y,t\Psi(y;\hat{s}_j)+(1-t)\Psi(y;\tilde{s}_j))dt
(\Psi(y;\hat{s}_j)-\Psi(y;\tilde{s}_j))
\end{eqnarray*}
where
$$
A_j(y)=\int_0^1\omega_p(y,t\Psi(y;\hat{s}_j)+(1-t)\Psi(y;\tilde{s}_j))dt.
$$
Since both $\Psi(y;\hat{s}_j)$ and $\Psi(y;\tilde{s}_j)$ approach $\Psi(y;s_\dag)$ when $j\to\infty$ on the interval $[M,0]$ for a fixed $M$, $M<y_j$, we have that
$$
|A_j(y)-\omega_p(y,\Psi(y;s_\dag))|\leq \varepsilon_j\;\;\mbox{on $[M,0]$},
$$
where $\varepsilon_j\to o$ as $j\to\infty$. This implies that $\Psi(y;\hat{s}_j)$ must coincide with $\Psi(y;\tilde{s}_j)$ for large $j$. This contradiction proves Lemma.

\end{proof}

\end{document}
